\section{Evaluation}
\label{sec:eval}

\subsection{Data Collection}
\label{sec:data}

\noindent\textbf{Human sessions.}
We recruited \textcolor{red}{XX} participants (age range \textcolor{red}{XX--XX},
\textcolor{red}{XX\%} male) across three naturalistic interaction tasks:
(T1)~social-media browsing (scrolling, tapping posts, entering search queries),
(T2)~mobile banking form filling (entering account numbers, amounts, confirming transactions),
and (T3)~prompted UI tapping (reacting to server-selected challenge buttons at random times).%
Participants held the phone in their preferred grip throughout; no instruction was given
regarding grip posture.%
Each session lasted 60\,s; we collected \textcolor{red}{XX} sessions per participant
across \textcolor{red}{XX} device models spanning four families:
Pixel (6/7/8/8a), Galaxy (S22/S23/A54), OnePlus (11/12), and Xiaomi (12/13)
(\textcolor{red}{XX} models total).%
All data was collected with IRB approval; participants provided informed written consent.%

\noindent\textbf{Agent sessions.}
Agent sessions were generated using three frameworks:
(i)~AutoGLM~\cite{autoglm2024} controlling an identical UI workflow via its default
ADB + cloud-VLM backend;
(ii)~AppAgent~\cite{appagent2024} performing the same tasks;
(iii)~custom ADB scripts replaying touch sequences without LLM inference.%
For each framework, we additionally tested A2-adversary (timing mimicry) sessions in
which the injection timestamps were resampled to match the lognormal distribution of
our human cohort.%
All agent sessions used the phone resting on a flat, non-vibrating surface (A0 default).%
A1-adversary sessions used an off-the-shelf vibration motor (ERM, 13\,mm diameter)
attached to the phone chassis and triggered by the ADB control host.%
A total of \textcolor{red}{XX} agent sessions were collected across frameworks and adversary levels.%

\noindent\textbf{Active-mode sessions.}
For active-mode evaluation, a separate split of \textcolor{red}{XX} human and
\textcolor{red}{XX} agent sessions included server-selected challenge events at
$t^* \in [\frac{T}{4}, \frac{3T}{4}]$ with a live challenge UI element displayed
at $t^*$.%

\subsection{Baselines}
\label{sec:baselines}

We compare PCCN against four baselines.%
\textbf{BeCAPTCHA-Accel}~\cite{becaptcha2021}: the accelerometer-only binary classifier
from BeCAPTCHA, retrained on our dataset.%
\textbf{TToS-Features}~\cite{ttos2026}: touch-kinematic features (coordinate centroid,
velocity profile entropy, inter-touch interval statistics) from Turing Test on Screen,
fed to a gradient-boosted tree classifier.%
\textbf{MiniRocket}~\cite{minirocket2021}: a state-of-the-art time-series classifier
applied to the raw accelerometer signal.%
\textbf{PCCN-Passive}: our model's passive head only (ablation).%

\subsection{Metrics}
\label{sec:metrics}

We report five metrics: overall accuracy (ACC), macro F1, AUC-ROC, equal error rate (EER),
false positive rate (FPR, human misclassified as agent), and false negative rate (FNR,
agent misclassified as human).%
EER is the standard metric for INFOCOM liveness detection papers.%
All results use a 5-fold cross-validation with a zero-overlap device-model split:
no device model in the test fold appears in the training fold.%

\subsection{Main Classification Results}
\label{sec:main}

Table~\ref{tab:main} reports passive-mode and active-mode results under the A0 (no evasion)
adversary condition.%

\begin{table}[t]
\centering
\caption{Classification results under A0 adversary. Passive mode: 2-class (Human vs.\ Agent).
Active mode: 3-class (Human / Autonomous / Overlay). All metrics are mean $\pm$ std
over 5 device-model folds.}
\label{tab:main}
\begin{tabular}{lccccc}
\toprule
Method & Mode & ACC & F1 & AUC & EER \\
\midrule
BeCAPTCHA-Accel~\cite{becaptcha2021} & Passive & \textcolor{red}{XX.X} & \textcolor{red}{XX.X} & \textcolor{red}{XX.X} & \textcolor{red}{XX.X} \\
TToS-Features~\cite{ttos2026}        & Passive & \textcolor{red}{XX.X} & \textcolor{red}{XX.X} & \textcolor{red}{XX.X} & \textcolor{red}{XX.X} \\
MiniRocket~\cite{minirocket2021}     & Passive & \textcolor{red}{XX.X} & \textcolor{red}{XX.X} & \textcolor{red}{XX.X} & \textcolor{red}{XX.X} \\
PCCN-Passive (ours)                  & Passive & \textcolor{red}{XX.X} & \textcolor{red}{XX.X} & \textcolor{red}{XX.X} & \textcolor{red}{XX.X} \\
\midrule
PCCN-Full (ours)                     & Active  & \textcolor{red}{XX.X} & \textcolor{red}{XX.X} & \textcolor{red}{XX.X} & \textcolor{red}{XX.X} \\
\bottomrule
\end{tabular}
\end{table}

\noindent\textbf{Key findings.}
PCCN-Full achieves \textcolor{red}{>XX\%} accuracy in active mode, outperforming all
passive baselines by \textcolor{red}{XX} percentage points.%
The kinematic baseline (TToS-Features) performs near chance on A2-timing-mimicry sessions
(see Section~\ref{sec:adversarial}), confirming the limitation identified in
Section~\ref{sec:bg_kinematic}.%
BeCAPTCHA-Accel achieves moderate accuracy on A0 sessions but degrades substantially
on A1 (vibration motor) sessions; its feature engineering does not capture the spectral
distinction between motor oscillation and finger-impact broadband transients.%

\subsection{Cross-Device Generalization}
\label{sec:crossdevice}

Fig.~\ref{fig:crossdevice} shows per-fold accuracy under the zero-overlap device-model
split, with training and test folds partitioned by manufacturer family.%
PCCN-Full achieves \textcolor{red}{>XX\%} across all four family splits, demonstrating
that the physical causality gap is device-agnostic.%
The smallest accuracy is observed when training on Pixel/OnePlus and testing on Xiaomi,
attributable to the higher baseline accelerometer noise floor on Xiaomi hardware;
PCCN compensates via the per-session z-score normalization in preprocessing.%

\begin{figure}[t]
  \centering
  \textcolor{green}{[Figure: Grouped bar chart. X-axis: test device family
  (Pixel, Galaxy, OnePlus, Xiaomi). Y-axis: accuracy (\%), 80--100 range.
  Four bar groups, each with five bars for: BeCAPTCHA-Accel (steel gray),
  TToS-Features (light blue), MiniRocket (green), PCCN-Passive (amber orange),
  PCCN-Full (navy blue). PCCN-Full bar consistently tallest.
  Error bars ($\pm$1 std). Dashed horizontal reference line at human-level benchmark.
  White background, unified palette.]}
  \caption{Cross-device generalization under a zero-overlap device-model split.
  PCCN-Full maintains $>$\textcolor{red}{XX\%} accuracy across all four device
  families, confirming device-agnostic operation.}
  \label{fig:crossdevice}
\end{figure}

\subsection{Adversarial Robustness}
\label{sec:adversarial}

Table~\ref{tab:adversarial} reports active-mode accuracy of PCCN-Full and baselines
under A0--A2 adversary conditions.%

\begin{table}[t]
\centering
\caption{Active-mode accuracy under adversary levels A0--A2.
Best result per adversary level in bold.}
\label{tab:adversarial}
\begin{tabular}{lcccc}
\toprule
Method & A0 (default) & A1 (vibration) & A2 (timing mimicry) \\
\midrule
BeCAPTCHA-Accel  & \textcolor{red}{XX.X} & \textcolor{red}{XX.X} & \textcolor{red}{XX.X} \\
TToS-Features    & \textcolor{red}{XX.X} & \textcolor{red}{XX.X} & \textcolor{red}{XX.X} \\
MiniRocket        & \textcolor{red}{XX.X} & \textcolor{red}{XX.X} & \textcolor{red}{XX.X} \\
PCCN-Full (ours) & \textbf{\textcolor{red}{XX.X}} & \textbf{\textcolor{red}{XX.X}} & \textbf{\textcolor{red}{XX.X}} \\
\bottomrule
\end{tabular}
\end{table}

\noindent\textbf{A1 (vibration motor).}
PCCN-Full retains \textcolor{red}{>XX\%} accuracy against A1 adversaries.%
The cross-modal attention layer learns to distinguish broadband finger-impact transients
from the narrowband sinusoidal oscillation produced by the ERM motor.%
The spectral centroid of the motor signal is concentrated at the motor's driving frequency
(\textcolor{red}{$\approx$XX}\,Hz), while genuine finger impact spans 1--50\,Hz with
an exponential spectral rolloff.%
The Transformer encoder captures this distinction via learned frequency-sensitive filters
in its self-attention heads.%

\noindent\textbf{A2 (timing mimicry).}
Timing mimicry reduces TToS-Features accuracy to near-chance but has no effect on
PCCN-Full, which relies on physical impulse presence rather than timing statistics.%
PCCN-Passive (which uses accelerometer context only, without timing stream) similarly
retains high accuracy under A2, confirming that the physical causality signal drives detection.%

\subsection{Overlay Agent Detection}
\label{sec:overlay}

Fig.~\ref{fig:overlay} shows the detection window sweep for overlay agent sessions as
a function of the number of challenge events $t^*$ observed in the session.%
With a single challenge event (\textit{n=1}), PCCN-Full achieves
\textcolor{red}{>XX\%} detection accuracy at FPR \textcolor{red}{$<$XX\%}.%
The key insight is that an overlay agent replaces exactly the challenged touch with an ADB
injection; the cross-modal attention head detects the impulse absence at that specific
token while genuine human impulses remain at all other positions.%

\begin{figure}[t]
  \centering
  \textcolor{green}{[Figure: Line plot. X-axis: number of challenge events $n$ (1, 2, 3, 4, 5).
  Y-axis: detection accuracy (\%), 60--100.
  Two lines: PCCN-Full (navy blue, rising from $\sim$XX\% at n=1 to $\sim$XX\% at n=5),
  PCCN-Passive (amber orange, flat near $\sim$XX\% regardless of n, since it cannot
  exploit active challenges).
  Shaded confidence bands. Dashed horizontal line at 95\% reference.
  White background, unified palette.]}
  \caption{Overlay agent detection accuracy vs.\ number of server-selected challenge
  events. A single challenge event is sufficient for $>$\textcolor{red}{XX\%} detection;
  passive mode cannot distinguish overlay agents from human sessions.}
  \label{fig:overlay}
\end{figure}

\subsection{Ablation Study}
\label{sec:ablation}

Table~\ref{tab:ablation} ablates the four components of PCCN-Full: the accelerometer
Transformer encoder, cross-modal attention (causal impulse detector), timing distribution
encoder, and Stage~1 MAE pre-training.%

\begin{table}[t]
\centering
\caption{Ablation study. Active-mode accuracy under A0. Each row removes one component.}
\label{tab:ablation}
\begin{tabular}{lcc}
\toprule
Configuration & ACC & EER \\
\midrule
PCCN-Full                             & \textcolor{red}{XX.X} & \textcolor{red}{XX.X} \\
w/o Cross-modal attention             & \textcolor{red}{XX.X} & \textcolor{red}{XX.X} \\
w/o Timing distribution encoder      & \textcolor{red}{XX.X} & \textcolor{red}{XX.X} \\
w/o MAE pre-training (Stage~1)        & \textcolor{red}{XX.X} & \textcolor{red}{XX.X} \\
Passive head only (no active mode)    & \textcolor{red}{XX.X} & \textcolor{red}{XX.X} \\
\bottomrule
\end{tabular}
\end{table}

Removing cross-modal attention causes the largest drop (\textcolor{red}{XX} pp),
confirming that explicit causal impulse detection is the primary signal.%
Removing the timing encoder causes a smaller drop (\textcolor{red}{XX} pp) that is
concentrated in the A2-timing-mimicry condition (where timing mimicry eliminates
the timing signal, leaving physical causality as the sole discriminant).%
MAE pre-training contributes \textcolor{red}{XX} pp, consistent with its role in
generalizing the encoder to unseen device models.%

\subsection{Attention Visualization}
\label{sec:attention}

Fig.~\ref{fig:attention} shows the cross-modal attention weights averaged over
\textcolor{red}{XX} human and \textcolor{red}{XX} agent test sessions, plotted
as a heatmap of attention weight vs.\ time offset relative to each touch event.%
Without explicit supervision on the [0,\,80]\,ms window, PCCN independently
concentrates attention at offsets 0--60\,ms post-touch for human sessions.%
Agent sessions show uniformly diffuse attention, consistent with the absence of
a localized impulse.%
This post-hoc interpretability confirms that PCCN has learned the physical causality
mechanism, not a spurious correlation.%

\begin{figure}[t]
  \centering
  \textcolor{green}{[Figure: Two-panel heatmap. Left panel (Human sessions):
  X-axis: time offset relative to touch event (ms), $-100$ to $+200$.
  Y-axis: touch event index (ordered by session time).
  Color: cross-modal attention weight (white$\rightarrow$amber, higher = more attention).
  Strong amber band at 0--60\,ms, consistent across all event rows.
  Right panel (Agent sessions): same axes.
  Color: uniformly pale (near-zero attention everywhere).
  Vertical dashed line at $t{=}0$ in both panels.
  Shared color bar on right.
  White background, unified palette.]}
  \caption{Cross-modal attention visualization. For human sessions (left), PCCN
  concentrates attention in the [0,\,60]\,ms post-touch window without explicit supervision,
  rediscovering the physical impulse response. Agent sessions (right) show diffuse attention,
  reflecting the absence of a causal impulse.}
  \label{fig:attention}
\end{figure}

\subsection{System Overhead}
\label{sec:overhead}

We measure client-side overhead on a Pixel 8 running Android 14.%
Sensor registration at \texttt{SENSOR\_DELAY\_FASTEST} and touch-event logging consume
\textcolor{red}{XX}\,mAh over a 30-s session (measured via Battery Historian),
equivalent to \textcolor{red}{XX\%} of idle power consumption.%
No background service is required; the agent runs within the host application process.%
Server-side PCCN inference on a 30-s window takes \textcolor{red}{XX}\,ms on a
single Intel Xeon core, well within the latency budget of any authentication flow.%
Peak memory consumption on the server is \textcolor{red}{XX}\,MB per concurrent session.%
